\chapter{Installation de Visual Studio Code}
\label{annexe:vscode}
Ouvrir un terminal et exécuter les commandes suivantes :

\begin{lstlisting}[style=commandstyle]
sudo apt update
sudo apt install curl apt-transport-https
\end{lstlisting}

Ajouter la clé GPG de Microsoft et le dépôt pour installer Visual Studio Code :

\begin{lstlisting}[style=commandstyle]
curl -sSL https://packages.microsoft.com/keys/microsoft.asc | sudo gpg --dearmor -o /usr/share/keyrings/packages.microsoft.gpg
\end{lstlisting}

Puis, ajouter le dépôt à la liste des sources du système :

\begin{lstlisting}[style=commandstyle]
echo "deb [arch=amd64 signed-by=/usr/share/keyrings/packages.microsoft.gpg] https://packages.microsoft.com/repos/code stable main" | sudo tee /etc/apt/sources.list.d/vscode.list > /dev/null
\end{lstlisting}

Mettre à jour les sources des paquets et installer Visual Studio Code :

\begin{lstlisting}[style=commandstyle]
sudo apt update
sudo apt install code
\end{lstlisting}

Si l'installation s'est bien déroulée, Visual Studio Code est disponible dans \textit{Show Applications} en bas à gauche sur Ubuntu. Il est recommandé de faire un clic droit sur l'application et de sélectionner \textit{Ajouter aux favoris} pour y accéder facilement.

Il est conseillé d'installer l'extension Python pour Visual Studio Code afin de bénéficier d'une prise en charge complète du langage, incluant l'auto-complétion et le débogage.
